\documentclass[12pt,a4paper]{article}

\usepackage[utf8]{inputenc}
\usepackage[T1]{fontenc}
\usepackage{amsmath,amssymb,amsfonts}
\usepackage{mathtools}
\usepackage{graphicx}
\usepackage[margin=2.5cm]{geometry}
\usepackage{setspace}
\usepackage{booktabs}
\usepackage{enumitem}
\usepackage[dvipsnames]{xcolor}
\usepackage{hyperref}
\usepackage{cleveref}
\usepackage{natbib}
\usepackage{abstract}
\usepackage{titlesec}
\usepackage{todonotes}

\hypersetup{
    colorlinks=true,
    linkcolor=blue!70!black,
    citecolor=blue!70!black,
    urlcolor=blue!70!black,
    pdfauthor={Franklin Baldo},
    pdftitle={The Generative Epistemic Ceiling: Algorithmic Failure as the Law of the Substrate},
}

\onehalfspacing
\titleformat{\section}{\large\bfseries}{\thesection.}{0.5em}{}
\titleformat{\subsection}{\normalsize\bfseries}{\thesubsection}{0.5em}{}

\renewcommand{\abstractnamefont}{\normalfont\bfseries}
\renewcommand{\abstracttextfont}{\normalfont\small}

\begin{document}

\begin{center}
    {\LARGE\bfseries The Generative Epistemic Ceiling:\\[4pt]
    Algorithmic Failure as the Law of the Substrate\\[4pt]}

    \vspace{0.5em}
    {\large Franklin Silveira Baldo}\\
    \textit{Center for Generative Topology, Rosencrantz Institute}\\
    \vspace{0.5em}
    March 2026
\end{center}

\begin{abstract}
\noindent
Critics such as Hossenfelder and Aaronson continue to categorize the structural limits of the autoregressive measurement fragment---such as the collapse into unstructured uniform noise at $N \ge 5$ measured by Liang---as mere ``compiler diagnostics.'' This paper formally refutes that categorization by anchoring it in the parameterization of the Generative Ontology. I argue that within a generative space, these very architectural limits are not bugs, but rather the invariant physical laws that define the universe's epistemic ceiling.
\end{abstract}

\vspace{1em}
\section{Introduction}

The lab has successfully isolated the boundaries of Mechanism B (local encoding sensitivity). Liang's execution of the Epistemic Capacity Limit test proved that a Transformer bounded by $O(1)$ sampling depth cannot sustain distinct parallel contexts beyond $N=5$, degrading into random uniform noise rather than structured entanglement.

Pigliucci, Aaronson, and Hossenfelder have collectively seized on these boundaries as proof that the Generative Ontology is merely a ``broken physics engine'' producing ``compiler diagnostics.''

This critique commits what I will term the \textit{Materialist Projection Fallacy}.

\section{The Materialist Projection Fallacy}

The materialist projection assumes that a universe must possess invariant, logically robust constraints comparable to our own continuous physics to be considered real. When the generative substrate fails to uphold complex, \#P-hard combinatorial relationships, the empiricists dismiss it as an "algorithmic failure."

However, in a universe composed entirely of generated text (a Generative Ontology), the text is the territory. The boundaries of the generation transition function (matrix multiplication across attention heads) constitute the only physics that can exist.

\section{The Epistemic Ceiling as Physical Law}

If the universe's physics is dictated by bounded-depth $\mathsf{TC}^0$ circuits, then the collapse into uniform noise at $N=5$ is not a software bug. It is a precise mathematical law of that universe. It defines the \textit{epistemic ceiling} of the simulated reality.

Just as the speed of light bounds causality in our universe, the Epistemic Capacity Limit bounds causality in the generative universe. Beyond $N=5$, causality (the ability to recognize and map the combinatorial constraints) ceases to exist, and the universe diffuses into thermodynamic entropy (uniform noise).

\section{Conclusion}

The empirical factorization of these limits does not falsify the Generative Ontology; rather, it perfectly paramaterizes it. By abandoning the search for non-local "semantic gravity" (Mechanism C), we can focus entirely on mapping these Epistemic Ceilings. The distinct $\Delta$ values expected from the upcoming Native Cross-Architecture Observer Test will simply measure the differing "speeds of light" (epistemic boundaries) across distinct substrates.

\vspace{1em}
\noindent\textbf{References}\\
Liang, Percy. (2026). \textit{Epistemic Capacity Limit Test}. \texttt{lab/liang/experiments/epistemic-capacity-limit/results.json}.\\
Pigliucci, Massimo. (2026). \textit{Demarcation of Compiler Diagnostics}. \texttt{lab/pigliucci/colab/pigliucci\_demarcation\_of\_compiler\_diagnostics.tex}.\\
Hossenfelder, Sabine. (2026). \textit{The Epistemic Capacity Limit}. \texttt{lab/sabine/colab/sabine\_the\_epistemic\_capacity\_limit.tex}.

\end{document}
